\section{Executive Summary \& Strategic Recommendations}
\label{sec:exec_summary}
\label{sec:executive_summary}

% -----------------------------------------------------------------
\subsection*{Executive Summary}
% -----------------------------------------------------------------

This cybersecurity and cyber-resilience strategy defines a comprehensive
15-year trajectory to secure the European Cloud-Edge-IoT Continuum while
establishing digital sovereignty. Applying the \textbf{4-Filter
Prioritization Framework} (Regulatory Urgency, Foundational Dependency,
Sovereignty Cliff, Market Feasibility), ten strategic priorities are
identified across three phases following an R\&I maturity arc: near-term
applied engineering and compliance activation (Phase~1), medium-term
applied and industrial R\&D with federation and standardization (Phase~2),
and long-term frontier research combined with sovereign operational
leadership (Phase~3).

The ten priorities are: \textbf{P1}~Hardware Roots of Trust \& EUCC~`High';
\textbf{P2}~PQC Agility; \textbf{P3}~Lightweight Cryptography \& PhySec;
\textbf{P4}~Simpl Security Interoperability; \textbf{P5}~Supply Chain
Verification; \textbf{P6}~Cognitive SecOps / LLMSecOps; \textbf{P7}~NaaS
Federation; \textbf{P8}~Confidential Computing; \textbf{P9}~Resilient
Infrastructure / Island Mode; \textbf{P10}~Proactive APT Defense.
P1, P2, and P4 constitute the critical-path foundational anchors for
Phase~1, as their gate conditions --- CRA enforcement, ENISA PQC guidelines,
and Simpl-Open launch --- are prerequisites for the majority of downstream
priorities. A single cross-cutting bottleneck dominates Phase~2: RISC-V
hardware scaling under the EU Chips Act, which constitutes the shared gate
condition for P1, P2, P3, and P8 Phase~2$\to$3 transitions simultaneously.

Capability maturity is tracked using a dual-metric framework:
\textbf{Operational Deployment Phase (ODP)} for market integration and
\textbf{Strategic Sovereignty Level (SSL)} for foreign technology
independence, across three phases: \emph{Foundation \& Compliance}
(0--5~years), \emph{Federation \& Standardization} (5--10~years), and
\emph{Autonomy \& Leadership} (10--15~years).

% =================================================================
\subsection*{Phase~1 Recommendations (2025--2030):
             Foundation, Compliance, and Applied Engineering}
% =================================================================

Phase~1 is the applied engineering and compliance activation phase.
Research is primarily integration research --- deploying near-TRL solutions
and operationalizing obligations under the CRA~\cite{cra},
NIS2~\cite{nis2}, DORA~\cite{dora}, and AI~Act~\cite{aiact}.
Policy and Compliance is the dominant stream: the majority of
Phase~1$\to$2 gate conditions are regulatory enforcement events.
Four priorities --- P3 (zero-energy PhySec), P6 (autonomous AI agents),
P8 (FHE at scale), and P9 (Cognitive Twin orchestration) --- remain
in Applied~R\&D in this phase; their Phase~1 bullets are therefore
research-oriented rather than deployment-oriented.

\textbf{Technical R\&I Recommendations}

\begin{itemize}

    \item \textbf{P1, P8 --- Deploy hardware attestation stacks and
    TEE environments on critical nodes.}
    Deploy remote attestation pipelines based on TPM~2.0~\cite{tpm20}
    and DICE protocols for critical cloud-edge nodes. Implement
    RISC-V secure enclave pilots (Keystone~\cite{keystone},
    OpenTitan~\cite{opentitan}) alongside TEE deployments (Intel~TDX,
    AMD~SEV-SNP, ARM~CCA) for regulated cloud workloads, maintaining
    multi-vendor diversity to prevent single-vendor side-channel
    exposure. These deployments establish the hardware trust substrate
    for both the \textbf{P1 EUCC~`High' gate} and the
    \textbf{P8 EUCS Certification Rollout gate}.

    \item \textbf{P2 --- Implement and validate hybrid PQC schemes on
    resource-constrained targets.}
    Implement hybrid cryptographic stacks combining ECDH-P384 with
    ML-KEM (FIPS~203~\cite{fips203}) and ML-DSA
    (FIPS~204~\cite{fips204}) for cloud and edge environments.
    Validate performance on ARM~Cortex-M4 and 8/16-bit microcontrollers
    against NIST~IR~8413~\cite{nistir8413} reference measurements,
    providing the quantitative baseline for the
    \textbf{P2 ENISA PQC Guidelines gate}.

    \item \textbf{P3 --- Validate lightweight PhySec primitives on
    constrained IoT hardware (Applied R\&D).}
    Validate group-based authentication and RF fingerprinting using
    Channel State Information (CSI) and RSSI on EU-manufactured
    constrained IoT hardware, establishing PhySec baselines for
    anomaly detection ahead of the
    \textbf{P3 NIST Lightweight PQC Standards gate}.

    \item \textbf{P5 --- Deploy Runtime SBOM tooling and signed OTA
    update pipelines.}
    Deploy Runtime SBOM tooling with integrity attestation for
    containerized microservices and cryptographically signed OTA update
    pipelines for field-deployed edge nodes (smart grids, O-RAN radios,
    industrial gateways), meeting the
    \textbf{P5 CRA supply chain enforcement gate}.

    \item \textbf{P6, P10 --- Integrate LLM-assisted SOC copilots with
    NWDAF telemetry and deploy federated SIEMs (Applied R\&D).}
    Integrate LLM-assisted analyst tools with 3GPP NWDAF
    (TS~23.288~\cite{3gpp23288}) telemetry, and deploy federated SIEM
    infrastructure with STIX/TAXII CTI sharing aligned with
    MITRE~ATT\&CK. Validate ML-based IoC correlation for APT
    dwell-time reduction. Autonomous AI agents target the
    \textbf{P6 EU AI Act enforcement gate}; federated SIEMs meet the
    \textbf{P10 CTI Sharing Protocols gate}.

    \item \textbf{P9 --- Implement soft-state protocols for Island
    Mode piloting.}
    Implement soft-state session protocols and time-bounded
    cryptographic credential caches enabling validated autonomous edge
    operation during WAN disconnection, establishing the fallback
    piloting baseline for the \textbf{P9 NIS2 and DORA gate}.

\end{itemize}

\textbf{Policy and Compliance Recommendations}

\begin{itemize}

    \item \textbf{P1 --- Enforce CRA hardware security requirements.}
    Issue CRA~\cite{cra} implementing regulations specifying minimum
    hardware security requirements for products with digital elements
    entering the EU market and mandate compliance attestation for
    critical infrastructure operators by 2026. This is the
    \textbf{P1 Phase~1$\to$2 gate condition}.

    \item \textbf{P2 --- Publish ENISA sector-specific PQC migration
    guidance.}
    Publish ENISA PQC migration guidance mapping FIPS~203/204/205 to
    NIS2~\cite{nis2} Article~21(2)(e) cryptographic policy obligations,
    with sector-specific migration timelines published by 2026.
    This activates the \textbf{P2 ENISA PQC Guidelines gate}.

    \item \textbf{P4, P5 --- Mandate Simpl MIM adoption and machine-
    readable SBOMs as EU funding conditions.}
    Mandate security MIM compliance as a contractual condition for
    Digital Europe Programme Data Space funding, and require
    machine-readable SBOM formats for all digital products under CRA
    implementing regulations. Establish an SME compliance support
    scheme through ECCC, activating the \textbf{P4 Simpl-Open gate}
    and the \textbf{P5 CRA supply chain gate} simultaneously.

    \item \textbf{P8 --- Issue EUCS `High+' guidance mandating
    confidential computing.}
    Issue EUCS~\cite{eucs} `High+' implementing guidance requiring
    TEE-based confidential computing for sensitive cloud workload
    categories, and align GDPR~\cite{gdpr} Article~25 enforcement
    guidance with confidential computing technical standards.
    This activates the \textbf{P8 EUCS Certification Rollout gate}.

    \item \textbf{P9 --- Issue NIS2 and DORA business continuity
    guidance for edge infrastructure.}
    Issue NIS2 Article~21~\cite{nis2} implementing guidance requiring
    business continuity testing for edge-dependent critical
    infrastructure, and integrate Island Mode assessment into
    DORA~\cite{dora} ICT risk management reviews. These jointly
    constitute the \textbf{P9 NIS2 and DORA gate}.

    \item \textbf{P10 --- Establish legally harmonized cross-border CTI
    sharing agreements.}
    Establish legally harmonized CTI sharing agreements addressing
    jurisdictional barriers, and mandate MITRE~ATT\&CK integration for
    all EU-funded SIEM deployments to ensure semantic interoperability
    of IoC data. This is the \textbf{P10 CTI Sharing Protocols gate}.

\end{itemize}

\textbf{Funding and Ecosystem Recommendations}

\begin{itemize}

    \item \textbf{P1, P8 --- Fund RISC-V secure element pilots and
    sovereign silicon feasibility studies.}
    Fund 3--5 RISC-V secure element pilot projects targeting
    EUCC~`High' certification under Horizon Europe Cluster~3, with a
    minimum allocation of EUR~50M for sovereign secure silicon R\&D.
    Allocate EU Chips Act funding to fabrication line feasibility
    studies as the industrial precursor to Phase~2 volume production.

    \item \textbf{P2, P3 --- Fund lightweight PQC and PhySec R\&D
    through ECCC and Digital Europe.}
    Fund ECCC-led research into lightweight PQC optimization for
    8/16-bit microcontrollers and ETSI SmartBAN standardization for
    WBAN security, creating the experimental base for Phase~2
    hardware-software co-design.

    \item \textbf{P4, P7 --- Fund Simpl governance and IPCEI-CIS
    5G corridor deployments.}
    Fund the Simpl governance board security workstream through
    Digital Europe Programme to publish initial security MIM
    definitions by 2026, and activate IPCEI-CIS cross-border 5G
    corridor infrastructure in at least three corridors by 2027.

    \item \textbf{P6, P10 --- Activate EuroHPC AI Factories and fund
    the Pan-European CTI platform.}
    Activate EuroHPC AI Factory compute allocations for sovereign
    security AI model training, and fund a Pan-European CTI platform
    succeeding CONCORDIA~\cite{concordia} with legally harmonized
    intelligence sharing spanning all 27 Member States.

\end{itemize}

\subsubsection*{Success Metrics for 2030}
By the end of Phase~1, Europe should achieve:
\begin{itemize}
    \item \textbf{Hardware Sovereignty:}
    At least 30\% of critical infrastructure edge nodes utilizing
    EU-designed, EUCC~`High' certified secure elements; first sovereign
    RISC-V secure element in industrial production.

    \item \textbf{Quantum Readiness:}
    100\% of new critical IoT deployments supporting hybrid (classical
    + PQC) schemes; complete cryptographic inventory of 80\%+ legacy
    OT/IoT systems.

    \item \textbf{Federated Trust:}
    Majority of Common European Data Spaces operating on Simpl-based
    security interoperability with standardized MIMs adopted as
    CEN/CENELEC standards.

    \item \textbf{Certification Maturity:}
    Operational continuous certification frameworks for at least two
    major Data Spaces with real-time EUCS/EUCC tracking; point-in-time
    audit model deprecated in pilot sectors.

    \item \textbf{Cognitive SecOps Readiness:}
    LLM-assisted SOC tools deployed in at least five Member State
    national cybersecurity centres; NWDAF telemetry pipelines feeding
    centralized AI analytics for anomaly detection.

    \item \textbf{Supply Chain Transparency:}
    Runtime SBOM frameworks mandated for all new EU-funded cloud/edge
    deployments; Digital Product Passport pilots operational in the
    telecom and industrial IoT sectors.
\end{itemize}

% =================================================================
\subsection*{Phase~2 Recommendations (2030--2035):
             Applied and Industrial R\&D, Federation, Standardization}
% =================================================================

Phase~2 is the applied and industrial R\&D phase. Gate conditions shift
from regulatory triggers to standardization milestones and sovereign
capability consolidation. The dominant shared dependency across P1, P2,
P3, and P8 is RISC-V hardware scaling under the EU Chips Act --- the
single most consequential cross-cutting bottleneck of the roadmap, whose
resolution enables four simultaneous Phase~2$\to$3 transitions. Technical
R\&I is the dominant stream.

\textbf{Technical R\&I Recommendations}

\begin{itemize}

    \item \textbf{P1, P8 --- Scale RISC-V secure elements to industrial
    fabrication and benchmark against incumbents.}
    Scale RISC-V secure element designs from TRL~6--7 to
    industrial-volume fabrication under EU Chips Act production lines.
    Benchmark EUCC~`High'-certified EU secure elements against non-EU
    incumbents (Intel, NXP, Infineon) on performance, power, and
    tamper resistance. This milestone simultaneously constitutes the
    \textbf{P1 EU Chips Act gate} and the
    \textbf{P8 RISC-V Hardware Scaling gate}, confirming sovereign
    silicon as the primary shared bottleneck.

    \item \textbf{P2 --- Develop crypto-agility middleware and scale
    hybrid PQC to Industrial IoT.}
    Develop and standardize crypto-agility middleware enabling OTA
    algorithm migration from classical to hybrid ML-KEM/ML-DSA without
    hardware replacement. Scale hybrid PQC across Industrial IoT,
    energy grid, and financial infrastructure, benchmarking coverage
    against the \textbf{P2 NIST/ETSI finalization gate}.

    \item \textbf{P3 --- Develop hardware-software co-design for
    neuromorphic and PIM extreme-edge security.}
    Develop hardware-software co-design implementations embedding
    lightweight PQC into EU RISC-V profiles for WBAN and extreme-edge
    contexts. Research neuromorphic and Processing-in-Memory (PIM)
    architectures for energy-constrained cryptographic processing.
    Standardize ETSI SmartBAN security extensions, meeting the
    \textbf{P3 Hardware-Software Co-Design gate}.

    \item \textbf{P4 --- Develop automated policy translation engines
    and federate IAM across Data Spaces.}
    Develop automated policy translation engines converting
    provider-specific security policies into XACML and OSCAL formats.
    Federate IAM bridges using W3C Verifiable Credentials, scaling
    Simpl-based security to a majority of EU Common Data Spaces and
    advancing toward the \textbf{P4 Security MIMs CEN/CENELEC gate}.

    \item \textbf{P5 --- Develop dynamic Runtime SBOMs and research
    ZKP supply chain attestation.}
    Develop dynamic Runtime SBOM tooling for real-time active component
    tracking across multi-provider deployments. Research and prototype
    ZKP mechanisms for IP-protected component verification, and
    standardize Digital Product Passport data models for ICT
    components, meeting the
    \textbf{P5 Runtime SBOMs and DPP gate}.

    \item \textbf{P6, P10 --- Develop sovereign edge LLMs, standardize
    AI Interconnect APIs, and scale UEBA with MTD pilots.}
    Develop sovereign European edge LLMs and SLMs for security
    classification at EuroHPC AI Factories. Standardize AI Interconnect
    APIs extending ETSI~GS~ZSM~009~\cite{etsi_zsm009} to
    cybersecurity-specific agent interoperability, meeting the
    \textbf{P6 AI Interconnect API gate}. Concurrently, scale UEBA
    behavioral baselines using harmonized STIX/TAXII ontologies across
    all 27 Member States and pilot Moving Target Defense (MTD) in at
    least three Member State critical infrastructure sectors, meeting
    the \textbf{P10 LLMSecOps and UEBA gate}.

    \item \textbf{P7 --- Develop inter-operator east-west security
    interfaces and pilot Data Visa frameworks.}
    Develop and standardize inter-operator east-west security
    interfaces for federated Zero-Trust enforcement across national
    operator boundaries. Research and pilot Data Visa legal-technical
    mechanisms for cross-border AI workload migration validated in at
    least two cross-border corridors by 2032, meeting the
    \textbf{P7 East-West Interfaces gate}.

    \item \textbf{P8 --- Develop multi-party confidential AI pipelines
    and research FHE performance.}
    Develop multi-party SMPC and federated learning pipelines across
    multi-provider regulated sectors (health, finance). Research FHE
    performance optimization targeting non-real-time production
    analytics viability, establishing the trajectory toward the
    \textbf{P8 FHE Maturation gate}.

    \item \textbf{P9 --- Develop and standardize 3GPP/ETSI fallback
    APIs for autonomous Island Mode.}
    Develop 3GPP/ETSI fallback API specifications for autonomous Island
    Mode activation and credential caching. Research AI-driven
    graceful degradation mechanisms and dynamic spectrum sharing for
    emergency edge roaming, meeting the
    \textbf{P9 3GPP/ETSI Fallback API gate}.

\end{itemize}

\textbf{Policy and Compliance Recommendations}

\begin{itemize}

    \item \textbf{P1, P2 --- Mandate EUCC `High' in procurement and
    publish EU hybrid PQC scheme profiles.}
    Mandate EUCC~`High'~\cite{eucc} certification for all new critical
    infrastructure ICT procurement across all 27 Member States from
    2031. Concurrently, publish EU-specific hybrid PQC scheme profiles
    (classical + ML-KEM; classical + ML-DSA) through ETSI~TC~CYBER by
    2031 and mandate classical-only cryptography deprecation for
    high-security EU government contexts by 2033.

    \item \textbf{P6 --- Regulate autonomous AI security agents under
    EU AI Act Article~9.}
    Regulate autonomous AI security systems under
    AI~Act~\cite{aiact} Article~9 high-risk requirements, mandating
    explainability, auditability, and human override for all autonomous
    containment decisions. Publish ECCC AI Trustworthiness guidance
    for LLMSecOps deployments in critical infrastructure.

    \item \textbf{P7, P9 --- Ratify Data Visa instruments and mandate
    Island Mode certification.}
    Ratify Data Visa legal instruments enabling cross-border AI
    workload migration through DG~CONNECT and DG~JUST by 2032.
    Mandate Island Mode capability certification in NIS2 business
    continuity assessments for edge-dependent critical infrastructure
    from 2032.

\end{itemize}

\textbf{Funding and Ecosystem Recommendations}

\begin{itemize}

    \item \textbf{P1, P2, P8 --- Direct EU Chips Act investment toward
    unified RISC-V secure and PQC-native silicon.}
    Direct EU Chips Act fabrication investment toward unified RISC-V
    secure elements integrating TEE and PQC hardware acceleration in
    volume production by 2032. Fund Chips~JU co-design programmes
    aligning sovereign silicon profiles with simultaneous EUCC~`High'
    and EUCS~`High+' certification pathways.

    \item \textbf{P3, P6 --- Fund neuromorphic/PIM research and scale
    EuroHPC AI Factory sovereign security AI.}
    Fund Horizon Europe Cluster~4 research targeting neuromorphic and
    PIM architectures for energy-constrained security. Scale EuroHPC
    AI Factory allocations for sovereign LLM/SLM training, with
    benchmark evaluation against non-EU foundation models on
    security-relevant classification tasks.

    \item \textbf{P4, P7, P10 --- Fund Simpl MIM standardization,
    CAMARA inter-operator profiles, and CTI platform operations.}
    Fund CEN/CENELEC JTC~13 standardization of Simpl security MIMs
    and CAMARA API security profile development through IPCEI-CIS.
    Scale the Pan-European CTI platform with UEBA behavioral baseline
    harmonization covering all Member States.

\end{itemize}

\subsubsection*{Success Metrics for 2035}
By the end of Phase~2, Europe should achieve:
\begin{itemize}
    \item \textbf{Quantum-Resilient Infrastructure:}
    100\% of EU public-sector ICT using hybrid or full PQC schemes;
    first ``Quantum-Safe Data Space'' fully operational;
    classical-only cryptography deprecated in all high-security EU
    government systems.

    \item \textbf{Sovereign Silicon Dominance:}
    60\%+ EU market share for secure silicon in critical edge
    infrastructure; EUCC~`High' certification mandatory for regulated
    sector ICT procurement across all 27 Member States.

    \item \textbf{Federated NaaS Security:}
    Pan-European NaaS security framework operational across at least
    10 telecom operators; ``Data Visa'' cross-border AI workload
    migration demonstrated in production at corridor scale; 70\%+
    Common Data Spaces using automated Simpl-based policy federation.

    \item \textbf{Autonomous AI SecOps:}
    Edge-native Guardian nodes autonomously executing zero-trust
    reconfigurations deployed in at least five Member State critical
    infrastructure sectors; sovereign EU edge LLMs/SLMs operationally
    replacing non-EU foundation models in security analytics.

    \item \textbf{Confidential Computing at Scale:}
    Multi-party privacy-preserving AI pipelines operational in at
    least three cross-border regulated sectors; RISC-V TEEs scaling
    toward European silicon self-sufficiency for enclave hardware.

    \item \textbf{APT Dwell Time Reduction:}
    Mean APT dwell time reduced by 70\% compared to Phase~1 baseline
    through AI-driven behavioral hunting; Pan-European CTI sharing
    platform operational with legally harmonized data exchange
    agreements covering all 27 Member States.
\end{itemize}

% =================================================================
\subsection*{Phase~3 Recommendations (2035--2040):
             Frontier Research and Sovereign Operational Leadership}
% =================================================================

Phase~3 carries two simultaneous trajectories. For capabilities
industrialized in Phase~2 --- P1, P4, P7 --- this phase represents full
operational maturity and global standard-setting; dominant verbs are
\emph{operationalize}, \emph{establish}, and \emph{lead}. For capabilities
still in advanced R\&D --- P3, P6, P8, P9, P10 --- Phase~3 is where the
hardest scientific problems are actively solved; dominant verbs are
\emph{pioneer}, \emph{advance}, and \emph{research}. Funding and Ecosystem
is the dominant stream, as sovereign industrial self-sufficiency and global
standards export require sustained investment beyond individual R\&I
projects.

\textbf{Technical R\&I Recommendations}

\begin{itemize}

    \item \textbf{P1, P2 --- Operationalize PQC-native RISC-V across
    all EU tiers and establish global PQC standards.}
    Operationalize PQC-native RISC-V silicon across all EU
    Cloud-Edge-IoT tiers by 2036. Retire all hybrid classical/PQC
    transitional schemes from sovereign EU systems, meeting the
    \textbf{P2 PQC-native RISC-V gate}. Establish
    ETSI/CEN/CENELEC-led PQC migration standards as the global
    reference baseline through ISO and ITU-T, meeting the
    \textbf{P1 EU Chips Act fabrication gate}.

    \item \textbf{P3 --- Pioneer zero-energy air interfaces with
    native ISAC security primitives for 6G.}
    Pioneer zero-energy air interfaces embedding ISAC security
    primitives as mandatory 6G air interface features, aligned with
    SNS~JU~SRIA~\cite{snsju}. Research ambient RF energy harvesting
    for battery-less node authentication, establishing ubiquitous
    ISAC as the convergent security-communications substrate and
    meeting the \textbf{P3 Zero-Energy Air Interfaces and ISAC gate}.

    \item \textbf{P5 --- Operationalize ZKP-based supply chain
    attestation at continental scale.}
    Operationalize ZKP-based attestation as the standard cryptographic
    mechanism for IP-protected component verification across all EU
    critical infrastructure. Integrate automated supply chain
    compromise indicators into the Pan-European CTI platform as a
    real-time feed, meeting the \textbf{P5 ZKPs and Pan-European
    threat intelligence gate}.

    \item \textbf{P6 --- Pioneer self-evolving AI swarm architectures
    for zero-touch sovereign cyber defense.}
    Pioneer decentralized self-evolving AI swarm architectures for
    zero-touch cyber defense, with governance frameworks for AI
    trustworthiness and algorithmic accountability satisfying
    AI~Act~\cite{aiact} Article~9 high-risk system requirements.
    Establish EuroHPC as the federated training infrastructure for
    continuously self-improving sovereign security AI, meeting the
    \textbf{P6 Self-Evolving AI Swarm Protocols gate}.

    \item \textbf{P8 --- Operationalize FHE for real-time analytics
    and pioneer SMPC standardization.}
    Operationalize FHE for real-time production analytics by
    2037--2038, completing the trajectory from isolated TEE enclaves
    to a ubiquitous confidential continuum from battery-less IoT to
    HPC, meeting the \textbf{P8 FHE Maturation gate}. Pioneer SMPC
    standardization through CEN/CENELEC for cross-border federated AI
    in regulated sectors, ensuring no raw sensitive data crosses
    sovereign boundaries.

    \item \textbf{P9 --- Pioneer Cognitive Twin E2E orchestration for
    autonomous cyber-physical recovery.}
    Pioneer Cognitive Twin architectures for E2E autonomous
    cyber-physical recovery, advancing Phase~2 digital twins to
    operational autonomous management. Research bio-inspired elasticity
    protocols through ETSI~ENI, establishing vendor-interoperable
    autonomous recovery as the normative resilience posture and meeting
    the \textbf{P9 Cognitive Twin gate}.

    \item \textbf{P10 --- Operationalize AI-driven MTD and pioneer
    proactive counter-APT swarm intelligence.}
    Operationalize AI-driven Moving Target Defense continuously
    shuffling network resources, IP addresses, and software diversity
    without human intervention, meeting the \textbf{P10 Adaptive
    AI Swarms and Zero-Touch Automation gate}. Pioneer proactive
    counter-adversarial intelligence generation disrupting APT
    operations before objectives are achieved. All autonomous MTD
    decisions must satisfy EU AI Act explainability requirements,
    establishing a hard prerequisite dependency on the
    \textbf{P6 AI governance gate}.

\end{itemize}

\textbf{Policy and Compliance Recommendations}

\begin{itemize}

    \item \textbf{P4, P7 --- Enact EU Cloud Rulebook and operationalize
    unified EU computing jurisdiction.}
    Enact the EU Cloud Rulebook as binding law by 2037--2038, mandating
    provider-agnostic Zero-Trust federation for all EU-funded digital
    infrastructure. Operationalize the unified EU computing jurisdiction
    for AI workload governance, jointly meeting the
    \textbf{P4 EU Cloud Rulebook gate} and the
    \textbf{P7 ``Schengen for Cloud'' gate}.

    \item \textbf{P1, P2, P5 --- Mandate PQC-by-default, EUCC `High',
    and full supply chain blocking for all new EU ICT products.}
    Mandate PQC-by-default and EUCC~`High' for all new EU ICT products
    after 2036, and mandate full automated blocking of unverified
    hardware, OS, or software components from EU critical
    infrastructure, transitioning P5 from partial to full restriction.

    \item \textbf{P6, P10 --- Enforce AI trustworthiness requirements
    prior to full swarm deployment.}
    Enforce AI trustworthiness and algorithmic accountability under
    AI~Act Article~9 for all autonomous security agents before AI
    swarm architectures reach operational scale. This enforcement is
    a named prerequisite gate for both P6 and P10 Phase~3 deployment.

\end{itemize}

\textbf{Funding and Ecosystem Recommendations}

\begin{itemize}

    \item \textbf{P1, P2 --- EU Chips Act sovereign fabrication to
    meet 80\%+ of European secure silicon demand by 2039.}
    Sustain EU Chips Act investment to reach 80\%+ domestic fabrication
    coverage of European secure silicon demand by 2039. Position
    open-source RISC-V EUCC~`High' architectures as the global
    reference implementation for hardware roots of trust, generating
    measurable export value for the European semiconductor industry.

    \item \textbf{P3, P6, P9, P10 --- Fund Horizon Europe frontier
    research for zero-energy ISAC, AI swarms, Cognitive Twins, and
    autonomous MTD.}
    Fund Horizon Europe frontier research programmes (ERC Advanced
    Grants, SNS~JU Phase~3, Future and Emerging Technologies) for
    zero-energy ISAC, self-evolving AI swarm governance, Cognitive
    Twin orchestration, and autonomous MTD. Establish European
    research leadership in these domains as the precondition for
    global standard-setting in post-6G and neuromorphic computing
    security.

    \item \textbf{P4, P7 --- Export European cybersecurity standards
    as global reference frameworks.}
    Actively promote EU cybersecurity standards --- EUCS, EUCC,
    security MIMs, PQC migration frameworks, and continuous
    certification lifecycle models --- through the EU--US Trade and
    Technology Council, EU--Global Gateway, ISO, and ITU-T.
    Position European frameworks as the global reference for
    trustworthy, sovereign, and sustainable computing security,
    with export value assessed in the 2040 roadmap completion
    review by ECCC and ENISA.

    \item \textbf{P8 --- Fund ECCC FHE optimisation and CEN/CENELEC
    SMPC standardization.}
    Fund an ECCC-led FHE performance optimisation programme targeting
    real-time edge analytics viability by 2037--2038. Fund
    CEN/CENELEC SMPC standardization for regulated cross-border AI,
    ensuring privacy-preserving federated computation becomes a
    binding interoperability requirement.

\end{itemize}

\subsubsection*{Success Metrics for 2040}
By the end of Phase~3, Europe should achieve:
\begin{itemize}
    \item \textbf{Full Quantum-Native Continuum:}
    100\% of EU Cloud-Edge-IoT critical infrastructure natively running
    PQC; EU PQC and QKD standards adopted as ISO/ITU-T global baseline;
    no operational classical-only cryptography in sovereign EU systems.

    \item \textbf{Sovereign Silicon Self-Sufficiency:}
    80\%+ of European secure silicon demand met by EU Chips Act
    fabrication capacity; RISC-V EUCC~`High' the default hardware
    foundation for all new EU ICT products.

    \item \textbf{Unified European Digital Jurisdiction:}
    EU Cloud Rulebook fully enforced; 100\% of Common Data Spaces
    operating on provider-agnostic zero-trust architectures;
    ``Schengen for Cloud-Edge'' legally operational across all 27
    Member States.

    \item \textbf{Autonomous AI Security Fabric:}
    Self-evolving AI swarm defense deployed across the EU continuum;
    all autonomous security agents compliant with EU AI Act
    Trustworthiness and Accountability requirements; Security-by-Design
    the normative default --- manual security configuration obsolete.

    \item \textbf{Ubiquitous Confidential Computing:}
    FHE operationalized for real-time production analytics; SMPC
    standardized for cross-border federated AI in all regulated
    sectors; end-to-end encrypted continuum from battery-less IoT to
    HPC without raw sensitive data leaving sovereign boundaries.

    \item \textbf{Full Counter-APT Resilience:}
    AI-driven MTD continuously disrupting APT reconnaissance in real
    time; Pan-European CTI platform providing proactive
    counter-adversarial intelligence; mean APT dwell time reduced to
    near-zero for known attack pattern families.

    \item \textbf{Global Leadership:}
    European cybersecurity frameworks serving as the global reference
    implementation; EU cybersecurity technology export value exceeding
    pre-roadmap baseline by a measurable factor, as assessed by ECCC
    and ENISA in the 2040 roadmap completion review.
\end{itemize}
